% Generated by scripts/assemble.py; corrections belong in source/corrections.json.
\examhead{127}{郑州大学 2024 年数学分析试卷}{121}
% p122-b01; source PDF page 122
\par\Needspace{5\baselineskip}
\blockanchor{p122-b01}
\majorquestion{一}(10 分) 求函数极限 $\lim_{x\to 0^{+}}\left[\frac{1}{x(1 + x)^{\frac{1}{x}}} -\frac{1}{xe}\right]$.\par
% p122-b02; source PDF page 122
\par\Needspace{5\baselineskip}
\blockanchor{p122-b02}
\majorquestion{二}(10 分) 求定积分 $\int_{0}^{2024} \frac{x}{e^{2024 - x} + e^x} \, \mathrm{d}x$.\par
% p122-b03; source PDF page 122
\par\Needspace{5\baselineskip}
\par\Needspace{6\baselineskip}
\blockanchor{p122-b03}
\majorquestion{三}(10 分) 设椭圆 $\frac{x^2}{4} + \frac{y^2}{9} = 1$ 在点 $A\left(1, \frac{3\sqrt{3}}{2}\right)$ 的切线交 $y$ 轴于点 $B, L$ 是从 $A$ 到 $B$ 的直线段, 计算\par
\nopagebreak[4]
% p122-b04; source PDF page 122
\blockanchor{p122-b04}
\displayformula{\int_ {L} \left(\frac {\sin y}{x + 1} - \sqrt {3} y\right) \mathrm{d} x + \left(\cos y \ln (1 + x) + 2 \sqrt {3} x - \sqrt {3}\right) \mathrm{d} y.}
% p122-b05; source PDF page 122
\par\Needspace{5\baselineskip}
\blockanchor{p122-b05}
\majorquestion{四}(10 分) 讨论 $\int_0^{+\infty} \frac{\sin x}{x^{p-1} + \frac{1}{x}} \, \mathrm{d}x (p \geq 0)$ 的条件收敛和绝对收敛性.\par
% p122-b06; source PDF page 122
\par\Needspace{5\baselineskip}
\blockanchor{p122-b06}
\majorquestion{五}(10 分) 求函数 $f(x) = \int_{0}^{x} \frac{\ln(1 + 2t)}{t} \, \mathrm{d}t$ 的麦克劳林级数展开式.\par
% p122-b07; source PDF page 122
\par\Needspace{5\baselineskip}
\blockanchor{p122-b07}
\majorquestion{六}(15 分) 讨论方程 $f(x) = -\frac{1}{2} (1 + e^{-1}) + \int_{-1}^{1} |x - t| e^{-t^2} \, \mathrm{d}t = 0$ 在 $[-1, 1]$ 上根的个数.\par
% p122-b08; source PDF page 122
\par\Needspace{5\baselineskip}
\blockanchor{p122-b08}
\majorquestion{七}(15 分) 设 $f(x, y) = \begin{cases} \frac{|x|^{\alpha}|y|^{\alpha}}{x^2 + y^2}, & x^2 + y^2 \neq 0; \\ 0, & x^2 + y^2 = 0. \end{cases}$ 证明:\par
% p122-b09; source PDF page 122
\blockanchor{p122-b09}
(1) $\alpha > 1$ 时, $f(x, y)$ 在 $(0, 0)$ 连续.\par
% p122-b10; source PDF page 122
\blockanchor{p122-b10}
(2) $\alpha > \frac{3}{2}$ 时, $f(x, y)$ 在 $(0, 0)$ 可微.\par
% p122-b11; source PDF page 122
\par\Needspace{5\baselineskip}
\blockanchor{p122-b11}
\majorquestion{八}(15 分) 证明: 函数项级数 $\sum_{n=1}^{\infty} \frac{x^{2}}{(1 + x^{2})^{n}}$ 在 $(- \infty, +\infty)$ 上点点收敛, 但并非一致收敛.\par
% p122-b12; source PDF page 122
\par\Needspace{5\baselineskip}
\blockanchor{p122-b12}
\majorquestion{九}(35 分) 证明题.\par
% p122-b13; source PDF page 122
\blockanchor{p122-b13}
(1) (5 分) 已知 $\lim_{x \to +\infty} [f(x) - (ax + b)] = 0$, 求 $a, b$.\par
% p122-b14; source PDF page 122
\blockanchor{p122-b14}
(2) (10 分) 设 $f(x)$ 在 $[a, +\infty)$ 上一致连续, $g(x) \in C([a, +\infty))$, 且 $\lim_{x \to +\infty} [f(x) - g(x)] = 0$. 证明: $g(x)$ 在 $[a, +\infty)$ 上一致连续.\par
% p122-b15; source PDF page 122
\blockanchor{p122-b15}
(3) (5 分) 用 (2) 的结论说明 $f(x) = \frac{1}{x} + \ln (1 + e^x)$ 在 $[1, +\infty)$ 上一致连续.\par
% p122-b16; source PDF page 122
\par\Needspace{6\baselineskip}
\blockanchor{p122-b16}
(4) (15 分) 设函数 $f(x)$ 在 $(- \infty, +\infty)$ 上可导, 且存在常数 $a, b, c, d (a < c)$ 使得\par
\nopagebreak[4]
% p122-b17; source PDF page 122
\blockanchor{p122-b17}
\displayformula{\lim _ {x \to - \infty} [ f (x) - (a x + b) ] = 0, \lim _ {x \to + \infty} [ f (x) - (c x + d) ] = 0.}
% p122-b18; source PDF page 122
\blockanchor{p122-b18}
证明：对任意的 $\lambda \in (a, c)$，存在 $\xi \in (-\infty, +\infty)$，使得 $f'(\xi) = \lambda$\par
% p122-b19; source PDF page 122
\par\Needspace{5\baselineskip}
\blockanchor{p122-b19}
\majorquestion{十}(20 分) 证明:\par
% p122-b20; source PDF page 122
\blockanchor{p122-b20}
(1) $\int_0^{+\infty}\frac{\sin ux}{x}\mathrm{d}x$ 在任何不包含 $u = 0$ 的闭区间 $[a,b]$ 上一致收敛.\par
% p122-b21; source PDF page 122
\blockanchor{p122-b21}
(2) 已知 $\int_0^{+\infty} \frac{\sin x}{x} \, \mathrm{d}x = \frac{\pi}{2}$, 求 $F(x) = \int_0^{+\infty} \frac{\sin^2 xt}{t^2} \, \mathrm{d}t (x > 0)$.\par
